% !TeX root = ../../Book2.tex
% 纲要标题：### 20.3 量子平面与扭曲代数
% 纲要位置：计划纲要.md，第 1710 行。

\section{量子平面与扭曲代数}
\label{b2:sec:2003}

把交换关系中的修正换成一个非零参数，可以得到另一类可计算的代数。下面始终固定 \(q\in k^\times\)。参数为一时恢复交换多项式环；其他参数是否为单位根，将直接影响中心的大小。

% 纲要要点（供目录核对）：
%
% * 量子平面的定义；从移动 jr 对字母产生扭曲因子，用预定基上的结合乘法证明独立性
% * 扭曲多项式环的实例；量子关系齐次而 Weyl 修正低次，取伴随分次时保留的信息不同
% * 中心与参数的关系；有限中心秩要求参数的精确阶，区别于 Weyl 中心随特征变化的机制
%

\subsection{量子平面的定义}
\label{b2:subsec:2003:01}

\begin{definition}\label{b2:def:quantum-plane}
设 \(k\) 为域，\(q\in k^\times\)。将自由结合代数 \(k\langle x,y\rangle\) 对关系 \(yx=qxy\) 生成的双边理想取商，所得 \(k\) 代数称为参数 \(q\) 的\term[liangzi-pingmian]{量子平面}[quantum plane]，记为
\[
k_q[x,y]=k\langle x,y\rangle/(yx-qxy).
\]
这里不要求 \(x,y\) 可逆。将它们也取逆得到的是另一种代数，不能与当前定义混同。
\end{definition}
\symidx{\(k_q[x,y]\)：参数 \(q\ne0\) 的量子平面}

\begin{theorem}\label{b2:thm:quantum-plane-normal-basis}
设 \(k\) 为域，\(q\in k^\times\)，\(k_q[x,y]=k\langle x,y\rangle/(yx-qxy)\) 为量子平面，则全部 \(x^iy^j\)、\(i,j\ge0\)，构成它的一组 \(k\) 基，且对所有非负整数 \(i,j,r,s\)，有
\begin{equation}\label{b2:eq:quantum-plane-product}
(x^iy^j)(x^ry^s)=q^{jr}x^{i+r}y^{j+s}.
\end{equation}
\end{theorem}
\begin{proof}
用 \(yx=qxy\) 移动所有错序字母，得到正规字的生成性。在乘积 \(x^iy^jx^ry^s\) 中，中间的每个 \(y\) 都要越过 \(r\) 个 \(x\)，总共交换 \(jr\) 次，每次产生一个因子 \(q\)。因而
\[
y^jx^r=q^{jr}x^ry^j,
\]
这给出了\eqref{b2:eq:quantum-plane-product}，但尚不能保证正规字独立。按照这个乘法取以 \(e_{ij}\)、\(i,j\ge0\) 为基的向量空间，令 \(e_{ij}\) 对应预期的正规字 \(x^iy^j\)，定义
\[
e_{ij}e_{rs}=q^{jr}e_{i+r,j+s}.
\]
乘三个基向量 \(e_{ij},e_{rs},e_{uv}\) 时，两种结合方式的系数指数分别为 \(jr+(j+s)u\) 与 \(su+j(r+u)\)，它们相同，所以乘法结合，单位为 \(e_{00}\)。元素 \(e_{10},e_{01}\) 满足量子平面的关系，故由\cref{b2:prop:algebra-relations-universal}得到到该代数的同态。正规字 \(x^iy^j\) 恰映到 \(e_{ij}\)，其独立性由后者为基得到。
\end{proof}

关系本身齐次，所以取 \(\deg x=\deg y=1\)，量子平面就是分次代数，第 \(n\) 次分量维数为 \(n+1\)。它与二元多项式环具有相同的分次维数，但乘法不同；例如 \(q\ne1\) 时，\(yx-xy=(q-1)xy\ne0\)。对这个分次所给的字长滤过取伴随分次，仍得到量子平面自身。Weyl 关系 \(\mathrm{d}\cdot x-x\cdot\mathrm{d}=1\) 则含低次修正，常数在最高次层消失，所以它的伴随分次是交换多项式环。量子关系中的 \(q\) 保留在同一次数内，取伴随分次不会消去这个扭曲。

\subsection{扭曲多项式环的实例}
\label{b2:subsec:2003:02}

令 \(\sigma:k[x]\to k[x]\)、\(\sigma\bigl(f(x)\bigr)=f(qx)\)。因为 \(q\ne0\)，它是自同构，逆为代入 \(q^{-1}x\)。由关系逐项计算，得到
\[
y\cdot f(x)=\sigma\bigl(f(x)\bigr)y,
\qquad y^r\cdot f(x)=\sigma^r\bigl(f(x)\bigr)y^r.
\]
所以每个元素可以唯一写为 \(\sum_j f_j(x)y^j\)，而乘法为
\[
\bigl(f(x)y^i\bigr)\bigl(g(x)y^j\bigr)
=f(x)\sigma^i\bigl(g(x)\bigr)y^{i+j}.
\]
这称为一个\term[niuqu-duoxiangshi-huan]{扭曲多项式环}[skew polynomial ring]，记为 \(k[x][y;\sigma]\)。在\secref{b2:sec:2005}中，会把它推广为允许导子修正的 Ore 扩张。

\begin{proposition}\label{b2:prop:quantum-plane-domain}
设 \(k\) 为域，\(q\in k^\times\)，\(k_q[x,y]=k\langle x,y\rangle/(yx-qxy)\) 为量子平面，则它没有零因子。若以正规式中 \(y\) 的最高指数为次数，任意两个非零元素乘积的次数等于两因子次数之和。
\end{proposition}
\begin{proof}
设两因子的最高项为 \(f_m(x)y^m\)、\(g_n(x)y^n\)。乘积的 \(y^{m+n}\) 系数为 \(f_m(x)\sigma^m\bigl(g_n(x)\bigr)\)。\(\sigma\) 是自同构，且 \(k[x]\) 是整环，所以这个系数非零。其他项的 \(y\) 次数更低，不能消去它，得到结论。
\end{proof}

例如 \(y(x^2+x+1)=(q^2x^2+qx+1)y\)。这里移动系数改变的是整段多项式的输入；不能只给整个多项式乘上同一个 \(q\)。若 \(q=-1\) 且特征不为二，偶次幂与 \(y\) 交换，奇次幂则改变符号。

\subsection{中心与参数的关系}
\label{b2:subsec:2003:03}

\begin{proposition}\label{b2:prop:quantum-plane-center}
设 \(k\) 为域，\(q\in k^\times\)，\(k_q[x,y]=k\langle x,y\rangle/(yx-qxy)\) 为量子平面。若 \(q\) 在 \(k^\times\) 中具有无限阶，则 \(Z\bigl(k_q[x,y]\bigr)=k\)。若 \(q\) 的阶为有限整数 \(\ell\ge1\)，则
\[
Z\bigl(k_q[x,y]\bigr)=k[x^\ell,y^\ell].
\]
后一情形中，量子平面作为中心上的模自由，基为 \(x^iy^j\)、\(0\le i,j<\ell\)，秩为 \(\ell^2\)。
\end{proposition}
\begin{proof}
由\eqref{b2:eq:quantum-plane-product}，正规式 \(P=\sum c_{ij}x^iy^j\) 与 \(x\) 交换恰好要求 \((q^j-1)c_{ij}=0\)，与 \(y\) 交换恰好要求 \((q^i-1)c_{ij}=0\)。正规单项式基使这些条件逐项成立而不能相互抵消。无限阶时只剩 \(i=j=0\)；有限阶时，\(q^r=1\) 恰好等价于 \(\ell\mid r\)，故要求 \(\ell\mid i,j\)，得到中心公式。

对一般指数写 \(i=a\ell+r,j=b\ell+s\)，其中 \(0\le r,s<\ell\)，再把中心因子提出，得到所列有限生成组。若它们有中心系数关系，展开这些系数后便得到不同的正规单项式之间的 \(k\) 线性关系，所有系数必须为零。因此它们确为中心上的基。
\end{proof}

这里的 \(\ell\) 必须是 \(q\) 的确切阶。仅有 \(q^\ell=1\) 可以保证 \(x^\ell,y^\ell\) 属于中心，却不能保证它们生成整个中心；例如 \(q=1\) 时虽然 \(q^2=1\)，中心仍为 \(k[x,y]\)，不是 \(k[x^2,y^2]\)。特别地，\(q=1\) 对应 \(\ell=1\)；\(q=-1\) 且特征不为二时，确切阶为二，中心为 \(k[x^2,y^2]\)，代数在中心上的秩为四。在特征二中，\(-1=1\)，同一写法则恢复交换多项式环。

\ProseParagraphBreak
量子平面中心的上述变化由 \(q\) 的乘法阶决定。\cref{b2:thm:weyl-center-simplicity}中 Weyl 代数的中心变化则由域特征决定：特征零时中心只有 \(k\)，特征 \(p>0\) 时中心为 \(k[x^p,\mathrm{d}^p]\)，在中心上的秩为 \(p^2\)。例如在特征零中取 \(q=-1\)，量子平面已在中心上有限，而 Weyl 代数仍有标量中心。两个有限秩公式都来自将正规单项式的指数按中心幂的指数取余，但使这些幂成为中心元的条件不同。
